\documentclass[10pt,a4paper]{amsart}
\usepackage{amsmath,amssymb,amsthm,graphicx}
% Retained preamble for the Stacks Project chapter snapshot in this directory.
% The chapter's own preamble.tex is a copy of the shared Stacks Project preamble; the block
% below is sliced MECHANICALLY, line for line, from it. Nothing was added, removed or
% redefined; the wrapper class, the xr-hyper/hyperref wiring and the externaldocument list
% are omitted because the assembler supplies the document class.
%
% lines taken from preamble.tex: 3, 5, 9, 11, 12, 156, 157, 158, ... (51 lines total)
%
% package coverage restored from the chapter's own preamble (2026-10-13): xy, xy-2cell

\usepackage[all]{xy}
\xyoption{2cell}
\UseAllTwocells

\newenvironment{reference}{\comment}{\endcomment}
\newenvironment{slogan}{\comment}{\endcomment}
\newenvironment{history}{\comment}{\endcomment}
\theoremstyle{plain}
\newtheorem{theorem}[subsection]{Theorem}
\newtheorem{proposition}[subsection]{Proposition}
\newtheorem{lemma}[subsection]{Lemma}
\theoremstyle{definition}
\newtheorem{definition}[subsection]{Definition}
\newtheorem{example}[subsection]{Example}
\newtheorem{exercise}[subsection]{Exercise}
\newtheorem{situation}[subsection]{Situation}
\theoremstyle{remark}
\newtheorem{remark}[subsection]{Remark}
\newtheorem{remarks}[subsection]{Remarks}
\numberwithin{equation}{subsection}
\def\lim{\mathop{\mathrm{lim}}\nolimits}
\def\colim{\mathop{\mathrm{colim}}\nolimits}
\def\Spec{\mathop{\mathrm{Spec}}}
\def\Hom{\mathop{\mathrm{Hom}}\nolimits}
\def\Ext{\mathop{\mathrm{Ext}}\nolimits}
\def\SheafHom{\mathop{\mathcal{H}\!\mathit{om}}\nolimits}
\def\SheafExt{\mathop{\mathcal{E}\!\mathit{xt}}\nolimits}
\def\Sch{\mathit{Sch}}
\def\Mor{\mathop{\mathrm{Mor}}\nolimits}
\def\Ob{\mathop{\mathrm{Ob}}\nolimits}
\def\Sh{\mathop{\mathit{Sh}}\nolimits}
\def\NL{\mathop{N\!L}\nolimits}
\def\CH{\mathop{\mathrm{CH}}\nolimits}
\def\proetale{{pro\text{-}\acute{e}tale}}
\def\etale{{\acute{e}tale}}
\def\QCoh{\mathit{QCoh}}
\def\Ker{\mathop{\mathrm{Ker}}}
\def\Im{\mathop{\mathrm{Im}}}
\def\Coker{\mathop{\mathrm{Coker}}}
\def\Coim{\mathop{\mathrm{Coim}}}
\DeclareMathSymbol{\boxtimes}{\mathbin}{AMSa}{"02}
\def\QCohstack{\mathcal{QC}\!\mathit{oh}}
\def\Cohstack{\mathcal{C}\!\mathit{oh}}
\def\Spacesstack{\mathcal{S}\!\mathit{paces}}
\def\Quotfunctor{\mathrm{Quot}}
\def\Hilbfunctor{\mathrm{Hilb}}
\def\Curvesstack{\mathcal{C}\!\mathit{urves}}
\def\Polarizedstack{\mathcal{P}\!\mathit{olarized}}
\def\Complexesstack{\mathcal{C}\!\mathit{omplexes}}
\def\Pic{\mathop{\mathrm{Pic}}\nolimits}
\def\Picardstack{\mathcal{P}\!\mathit{ic}}
\def\Picardfunctor{\mathrm{Pic}}
\def\Deformationcategory{\mathcal{D}\!\mathit{ef}}
\usepackage{cleveref}
\crefname{originaldefinition}{original definition}{original definitions}
\Crefname{originaldefinition}{Original definition}{Original definitions}
\makeatletter
\newlabel{groupoids-definition-equivalence-relation}{{022P}{source}{}{sourcepointer-1}{}}
\newlabel{groupoids-definition-equivalence-relation@cref}{{[originaldefinition][0][]022P}{[1][0][]source}{}{}{}}
\makeatother
\crefname{originaltheorem}{original theorem}{original theorems}
\Crefname{originaltheorem}{Original theorem}{Original theorems}
\makeatletter
\newlabel{theorem-presentation}{{02WW}{source}{}{sourcepointer-2}{}}
\newlabel{theorem-presentation@cref}{{[originaltheorem][0][]02WW}{[1][0][]source}{}{}{}}
\makeatother
\title{A group action free on residue fields defines an etale equivalence relation}
\author{The Stacks Project Authors}
\date{stacks-project:spaces.tex:local-snapshot-20261012; GNU FDL 1.2 or later}
\begin{document}
\maketitle
\setcounter{subsection}{1}
\noindent\textit{Editorial note.} All delivered mathematical text is the author's own, copied from the pinned Stacks Project chapter file spaces.tex; no statement or proof step was rewritten, reordered or removed.\par\medskip
\noindent\textit{Editorial note.} The statement displays the map j twice (statement and first line of the proof); both displays are part of the counted spans and are preserved verbatim.\par\medskip
\noindent\textit{Editorial note.} The two references of the counted spans are recorded as non-load-bearing navigation pointers; the argument that consumes them stays inside the counted proof span.\par\medskip
\noindent\textit{Editorial note.} Preamble: sources/stacks-spaces/retained-preamble.tex, sliced mechanically from the chapter's own preamble.tex.\par\medskip
\noindent\textit{Allowed standard prerequisites (editorial).}\begin{itemize}
\item Definition of an equivalence relation on a category fibred in groupoids. Groupoids, Definition groupoids-definition-equivalence-relation, used at proof lines 2511-2513. Supplies the axioms verified at the end of the proof.
\item The presentation theorem for algebraic spaces. Theorem theorem-presentation, used at proof lines 2482-2486. Turns an etale equivalence relation into an algebraic space.
\end{itemize}\medskip
\ifdefined\hypertarget\hypertarget{sourcepointer-1}{}\fi
\noindent\textit{Source reference.} Original definition 022P: focused-seeds/AGP/sources/stacks-snapshot-20261011/tags.txt, source lines 2714--2714. Navigation pointer into the Stacks Project; the standard fact it supplies is applied inside the counted proof span and the argument resting on it stays there.\par\medskip
\ifdefined\hypertarget\hypertarget{sourcepointer-2}{}\fi
\noindent\textit{Source reference.} Original theorem 02WW: focused-seeds/AGP/sources/stacks-snapshot-20261011/tags.txt, source lines 3819--3819. Navigation pointer into the Stacks Project; the standard fact it supplies is applied inside the counted proof span and the argument resting on it stays there.\par\medskip
\begin{lemma}
\label{lemma-quotient}
Let $U \to S$ be a morphism of $\Sch_{fppf}$.
Let $G$ be an abstract group. Let $G \to \text{Aut}_S(U)$
be a group homomorphism. Assume
\begin{itemize}
\item[(*)] if $u \in U$ is a point, and $g(u) = u$
for some non-identity element $g \in G$, then $g$
induces a nontrivial automorphism of $\kappa(u)$.
\end{itemize}
Then
$$
j :
R = \coprod\nolimits_{g \in G} U
\longrightarrow
U \times_S U,
\quad
(g, x) \longmapsto (g(x), x)
$$
is an \'etale equivalence relation and hence
$$
F = U/R
$$
is an algebraic space by Theorem \ref{theorem-presentation}.
\end{lemma}

\begin{proof}
In the statement of the lemma the symbol $\text{Aut}_S(U)$ denotes
the group of automorphisms of $U$ over $S$.
Assume $(*)$ holds. Let us show that
$$
j :
R = \coprod\nolimits_{g \in G} U
\longrightarrow
U \times_S U,
\quad
(g, x) \longmapsto (g(x), x)
$$
is a monomorphism. This signifies that if $T$ is a nonempty
scheme, and $h : T \to U$ is a $T$-valued point such that
$g \circ h = g' \circ h$ then $g = g'$. Suppose
$T \not = \emptyset$, $h : T \to U$ and $g \circ h = g' \circ h$.
Let $t \in T$. Consider the composition
$\Spec(\kappa(t)) \to \Spec(\kappa(h(t))) \to U$.
Then we conclude that $g^{-1} \circ g'$ fixes $u = h(t)$ and
acts as the identity on its residue field. Hence $g = g'$ by $(*)$.

\medskip\noindent
Thus if $(*)$ holds we see that $j$ is a relation (see
Groupoids, Definition \ref{groupoids-definition-equivalence-relation}).
Moreover, it is an equivalence relation since on $T$-valued points
for a connected scheme $T$ we see that
$R(T) = G \times U(T) \to U(T) \times U(T)$ (recall that we always
work over $S$). Moreover, the morphisms $s, t : R \to U$ are \'etale
since $R$ is a disjoint product of copies of $U$.
This proves that $j : R \to U \times_S U$ is an \'etale equivalence relation.
\end{proof}

\end{document}
